% ============================================================================
% EXERCISE A: FALLOW WATER BALANCE
% ============================================================================
% Duration: 90 minutes (hands-on; Module 2 support)
% Learning Goals: Understand daily water accounting; sowing readiness criteria
% Prerequisites: Module 1, Module 2

\chapter{Exercise A: Fallow Water Balance}
\label{ch:exercise-a}

\section{Overview}

\textbf{Objective:} Analyze 7 years of fallow (no crop) simulation to understand water dynamics and identify sowing opportunities.

\textbf{File:} \filename{Exercise\_A\_Fallow\_Water\_Balance.apsimx}

\textbf{Time:} 90 minutes

\textbf{Deliverable:} Excel workbook with analysis, plots, and written interpretation.

\section{Simulation Setup}

\subsection{Configuration}

\begin{table}[H]
\centering
\caption{Exercise A Simulation Parameters}
\label{tab:exercise-a-setup}
\begin{tabularx}{\textwidth}{l|X}
\toprule
\textbf{Parameter} & \textbf{Value} \\
\midrule
Location & Northam, Western Australia (WA Wheatbelt) \\
Climate & 1950--1957 (7-year subset of 101-year record) \\
Weather file & AU\_Northam.met \\
Soil & Norwin Vertosol (Black Clay) \\
Crop & None (fallow) \\
Management & None \\
Reported variables & Clock.Today, Soil.ESW, Weather.Rainfall, Soil.Evaporation, Soil.Drainage \\
\bottomrule
\end{tabularx}
\end{table}

\subsection{Key Parameters}

\begin{description}
  \item[\textbf{ESW threshold}] 100 mm (typical sowing readiness limit)
  \item[\textbf{Rainfall threshold}] 25 mm in last 7 days (recent moisture indicator)
  \item[\textbf{Sowing window}] May 1 -- July 10 (typical Queensland wheat season)
\end{description}

\section{Task 1: Run Simulation}

\begin{enumerate}
  \item Open APSIM
  \item File $\to$ Open $\to$ \filename{Exercise\_A\_Fallow\_Water\_Balance.apsimx}
  \item Verify Clock is set to 1950-01-01 to 1957-12-31 (7 years)
  \item Click Simulation in tree
  \item Press Run (or Ctrl+R)
  \item Wait for completion (should be fast; $<$ 10 seconds)
  \item Right-click DataStore $\to$ Export $\to$ To Excel
  \item Save as \filename{Exercise\_A\_Results.xlsx}
\end{enumerate}

\section{Task 2: Data Exploration}

\subsection{Examination Checklist}

\begin{itemize}
  \item How many rows in the output? (Should be $\approx$ 2555 daily records for 7 years)
  \item What are the column names? (Should match Report variables: Date, ESW, Rainfall, Evaporation, Drainage)
  \item Any missing values? (Should be none; every day should have data)
  \item Date range correct? (1950-01-01 to 1957-12-31)
\end{itemize}

\subsection{Quick Statistics}

In Excel, calculate:
\begin{itemize}
  \item Average ESW across all days
  \item Maximum ESW per year
  \item Minimum ESW per year
  \item Total annual rainfall by year
\end{itemize}

\section{Task 3: Identify Sowing Opportunities}

\subsection{Criteria}

Sowing is possible when BOTH conditions are met:
\begin{itemize}
  \item ESW $>$ 100 mm
  \item Rainfall in last 7 days $>$ 25 mm
  \item Date between May 1 and July 10
\end{itemize}

\subsection{Excel Implementation}

\begin{enumerate}
  \item Create helper columns:
    \begin{itemize}
      \item Column F: 7-day cumulative rainfall (sum of Rainfall over current + previous 6 days)
      \item Column G: Date in sowing window? (IF(AND(Date$\geq$May1, Date$\leq$July10), 1, 0))
      \item Column H: ESW sufficient? (IF(ESW$>$100, 1, 0))
      \item Column I: Rainfall sufficient? (IF(F$>$25, 1, 0))
      \item Column J: Sowing ready? (IF(AND(G, H, I), 1, 0))
    \end{itemize}
  
  \item Count: How many days are sowing-ready? (COUNTIF(J, 1))
\end{enumerate}

\section{Task 4: Visualization}

\subsection{Plot 1: ESW Time Series}

\begin{itemize}
  \item X-axis: Date (all 7 years)
  \item Y-axis: ESW (mm)
  \item Add reference line: ESW = 100 mm (sowing threshold)
  \item Title: ``Extractable Soil Water: 7-Year Fallow''
\end{itemize}

\textbf{Interpretation:} When does ESW peak? When is it lowest? Pattern seasonality?

\subsection{Plot 2: Annual ESW Pattern}

\begin{itemize}
  \item Calculate monthly average ESW (group by year-month)
  \item Plot 7 overlaid lines, one per year (different colors)
  \item X-axis: Month (1--12)
  \item Y-axis: Average ESW (mm)
  \item Title: ``Monthly ESW Pattern: 7 Years''
\end{itemize}

\textbf{Interpretation:} Is seasonal pattern consistent across years? When does ESW typically peak?

\subsection{Plot 3: Rainfall vs. ESW Change}

\begin{itemize}
  \item Create column: Daily ESW change = ESW(today) - ESW(yesterday)
  \item Scatter plot: X = Daily Rainfall, Y = Daily ESW change
  \item Add trend line
  \item Title: ``Daily Rainfall vs. ESW Change''
\end{itemize}

\textbf{Interpretation:} How much does ESW change per mm of rainfall? Any threshold effects?

\section{Task 5: Analysis Questions}

Answer the following using your data:

\begin{enumerate}
  \item \textbf{Seasonal Pattern:} ESW is typically highest in [season?] and lowest in [season?]. Why?
  
  \item \textbf{Sowing Ready:} In how many days per year (on average) are conditions suitable for sowing (ESW $>$ 100 mm, rain threshold, date window)? What is the range across the 7 years?
  
  \item \textbf{Dry Year vs. Wet Year:} Identify the driest and wettest year in your data (by total rainfall). Compare their ESW profiles. What are the implications for management?
  
  \item \textbf{Water Loss Processes:} For a typical year, calculate total water loss to:
    \begin{itemize}
      \item Evaporation (sum of daily Evaporation values)
      \item Drainage (sum of daily Drainage values)
    \end{itemize}
    Which process dominates? How does it vary seasonally?
  
  \item \textbf{Sowing Decision Rule:} Propose a sowing rule for farmers in this region. Should they wait for ESW $>$ 100 mm? What rainfall accumulation period is most predictive (7 days? 14 days)? What date window makes sense?
\end{enumerate}

\section{Task 6: Comparison Scenario (Extension)}

\textit{Optional:} Modify the simulation to explore:

\subsection{Scenario A: Shallow Soil}

Create a variant with half the soil depth (lower water-holding capacity). How does ESW pattern change? Sowing opportunities more or less frequent?

\subsection{Scenario B: Different Climate}

Substitute a wetter region's weather file (or a drier region). Repeat analysis. Compare.

\section{Deliverable}

Submit a single \filename{Exercise\_A\_Report.xlsx} or \filename{Exercise\_A\_Report.docx} containing:

\begin{itemize}
  \item $\checkmark$ Analyzed data (with helper columns for sowing readiness)
  \item $\checkmark$ 3+ visualizations (plots as described above)
  \item $\checkmark$ Written answers to Analysis Questions (5 items)
  \item $\checkmark$ Brief reflection (100--150 words): What surprised you? How would you use this in a real farming decision?
\end{itemize}

\section{Grading Rubric}

\begin{table}[H]
\centering
\begin{tabularx}{\textwidth}{l|r|X}
\toprule
\textbf{Criterion} & \textbf{Points} & \textbf{Description} \\
\midrule
Simulation setup & 10 & File opened, run successfully, data exported \\
Data exploration & 10 & Correct record count, column names, date range verified \\
Sowing identification & 15 & Correct criteria applied; counts accurate \\
Visualization (3 plots) & 30 & Clear labels, interpretable scales, titles \\
Analysis answers & 25 & Thoughtful interpretation of data; connections to concepts \\
Reflection & 10 & Personalized insights; practical relevance \\
\midrule
\textbf{Total} & \textbf{100} & \\
\bottomrule
\end{tabularx}
\end{table}

% ============================================================================
% END EXERCISE A
% ============================================================================
